\section*{Chapitre \ref{ch:s2:delta-one-and-dividends} --- Delta-One et dividendes}

\begin{solution}{exo:s2:delta-one-and-dividends:1}
$(7 - 5)$ points de base sur 50~millions de dollars~: 10\,000~\$.
\end{solution}

\begin{solution}{exo:s2:delta-one-and-dividends:2}
$(40 - 15) \times 0.25 - 5 = 1.25$ points de base sur 100~millions de dollars~: 12\,500~\$ pour le trimestre.
\end{solution}

\begin{solution}{exo:s2:delta-one-and-dividends:3}
$108.7 / 95.7 - 1 = 13.6~\%$ au total, $1.136^{1/3} - 1 = 4.3~\%$ par an. La moyenne simulée, 3.8~\% par an, est plus faible parce que les années de récession réduisent
les dividendes, et la moyenne des rendements annualisés est inférieure au rendement annualisé moyen.
\end{solution}

\begin{solution}{exo:s2:delta-one-and-dividends:4}
Le souscripteur de la note est vendeur d'un put sur l'indice à travers la protection~; l'émetteur en est acheteur et se couvre en détenant l'indice. Des dividendes plus élevés abaissent le
forward de l'indice et augmentent la valeur du put pour l'émetteur, ce qui diminue la note qu'il doit~: une hausse de 0.1 point du rendement de dividende lui rapporte 0.11 pour 100.
\end{solution}

\begin{solution}{exo:s2:delta-one-and-dividends:5}
Si les émetteurs de produits structurés vendent des dividendes à court terme pour se couvrir, les acheteurs qui les prennent sont rémunérés par une décote, ce qui se traduit par un rendement espéré élevé
sur les créances de dividendes à court terme~; une prime de risque et une prime d'offre sont difficiles à distinguer à partir des seuls prix.
\end{solution}

\begin{solution}{exo:s2:delta-one-and-dividends:6}
Les banques sont déclarées et facturées sur leur bilan à la fin des trimestres, si bien qu'elles le réduisent alors~; les rares qui peuvent détenir le panier face au contrat à terme
exigent davantage pour cela, et le taux de financement implicite du contrat à terme augmente.
\end{solution}

\begin{solution}{exo:s2:delta-one-and-dividends:7}
Le P\&L moyen par transaction vaut 1.14 points de base à 1 minute, 0.77 à 2, 0.40 à 3, 0.05 à 4 et $-0.26$ à 5~: la transaction cesse de rapporter entre quatre et cinq
minutes. Avec une demi-vie de 10 minutes, chaque minute de retard fait perdre une partie d'un écart qui n'excède les coûts que de quelques points de base.
\end{solution}

\begin{solution}{exo:s2:delta-one-and-dividends:8}
Le rendement s'accompagne d'un risque actions~: la transaction à un an est corrélée à 0.41 avec l'indice, perd dans environ un cinquième des années et recule de 24.4~\% en moyenne dans une
année de récession, quand les obligations se comportent généralement bien. Elle est rémunérée pour fournir une assurance aux émetteurs de produits structurés, et non pour prêter.
\end{solution}

\begin{solution}{pb:s2:delta-one-and-dividends:1}
\begin{enumerate}
\item Négocier un contrat à terme sur indice contre son panier lorsqu'il s'écarte de sa fair value d'un montant supérieur au coût des deux, et déboucler la position lorsque l'écart se referme.
\item Écarts de valorisation à la minute, avec un écart-type de 4 points de base et une demi-vie de 10 minutes~; 2 points de base par côté pour le panier et 0.5 pour le
contrat à terme, 5 pour un aller-retour.
\item 2\,230 transactions par an~: 1.64 points de base chacune sans délai, 1.14 à une minute, 0.77 à deux, $-0.26$ à cinq.
\item L'entrée se fait au-delà du coût d'un aller-retour et la sortie à zéro, si bien que chaque transaction rapporte au moins l'excédent.
\item Détenir un actif face à un dérivé qui le livre plus tard, en gagnant le spread de financement implicite diminué du coût du financement et du capital.
\item Le prix au-dessus du comptant, diminué des dividendes, implique le taux auquel le marché finance le panier.
\item Au moment du roll~: 10~\% des trimestres, 15.4 points de base bloqués, 0.4 par an~; dans la fenêtre de fin de trimestre~: 67.5~\%, 40.7, 8.0 par an.
\item La banque détient le panier et verse au client son rendement, moyennant un spread au-dessus du financement.
\item Le souscripteur de la note est vendeur d'un put~; la couverture de l'émetteur le rend acheteur de dividendes~: 0.11 pour 100 pour 0.1 point de rendement, soit environ 1.1 million par milliard.
\item La décote des contrats à terme sur dividendes par rapport aux dividendes espérés apparaît lorsque les détenteurs naturels de dividendes vendent davantage que les autres ne sont prêts à acheter à la fair value.
\item 100~EUR par point des dividendes de l'EURO STOXX 50 sur l'année du contrat, réglé en espèces, jusqu'à dix ans d'échéance.
\item Contrats à terme de 96.7 à 93.7, face à des dividendes espérés de 100.8 à 117.2~; rendements moyens de 3.8--4.2~\% par an~; écart-type de la transaction à un an
12.3~\%, quantile à 5~\% $-25.5~\%$, probabilité de perte de 20.9~\%.
\item Conservé jusqu'à l'échéance, le contrat à terme sur dividendes à un an rapporte en moyenne 4.2~\% par an, grâce à une décote liée à l'offre de 4~\%, avec une perte dans environ une année sur cinq
et de $-24.4~\%$ une année de récession.
\item À partir de 2004, les swaps de dividendes ont été une source majeure de profit pour les hedge funds, car les dividendes implicites se négociaient avec de fortes décotes et les dividendes augmentaient.
\item Les créances de dividendes à court terme affichent des rendements espérés, des ratios de Sharpe et des volatilités plus élevés que l'indice, avec des bêtas inférieurs à un.
\item Un modèle factoriel à deux facteurs ajuste la courbe des contrats à terme sur dividendes, et les investisseurs ne révisent que peu les dividendes au-delà du cycle économique.
\item En récession, lorsque les dividendes sont réduits et que les actions baissent ensemble.
\item Financement par total return swap.
\item Les deux rapportent un spread en contrepartie de la mise à disposition de bilan, et les deux s'élargissent lorsque les bilans bancaires sont rares.
\item Capacité de bilan et de couverture~: le rendement de l'indice sans le financer, et le risque de dividende sans celui de l'indice.
\end{enumerate}
\end{solution}

\begin{solution}{iq:s2:delta-one-and-dividends:1}
Le comptant diminué de la valeur actuelle des dividendes jusqu'à l'échéance, capitalisé au taux de financement jusqu'à l'échéance. Les prévisions de dividendes, le taux de financement qu'utilise le marché
(et non le taux sans risque), la retenue à la source et le calendrier des dates de détachement peuvent toutes se révéler fausses.
\end{solution}

\begin{solution}{iq:s2:delta-one-and-dividends:2}
Par le bas~: les dividendes annoncés et prévus de chaque composante, les politiques de distribution et les bénéfices, pondérés par les poids de l'indice et les dates de détachement~; vérifiés
par rapport aux contrats à terme sur dividendes et aux forwards implicites des options.
\end{solution}

\begin{solution}{iq:s2:delta-one-and-dividends:3}
À partir du coût de détention du panier sur la période~: le financement de la banque et la charge de bilan, plus élevée en fin de trimestre, auxquels
s'ajoutent les dividendes reversés et le coût de négociation du panier~; comparer avec le financement implicite du contrat à terme sur les mêmes dates.
\end{solution}

\begin{solution}{iq:s2:delta-one-and-dividends:4}
Une récession avec des réductions de dividendes de 30--50~\% dans les années proches, corrélée à une baisse des actions, et une interdiction réglementaire des dividendes bancaires~; appels de marge sur
les contrats à terme à mesure que les cours baissent.
\end{solution}

\begin{solution}{iq:s2:delta-one-and-dividends:5}
Un flux de poids d'indice, d'opérations sur titres et de prévisions de dividendes~; des courbes de financement~; le comptant du panier en temps réel~; un recalcul à chaque modification d'une entrée~;
des alertes lorsque le contrat à terme s'écarte de la fair value~; et des entrées versionnées pour l'audit.
\end{solution}

\begin{solution}{iq:s2:delta-one-and-dividends:6}
$F(D + \delta) - F(D) = -\delta e^{rT}$. Un trader qui sous-estime les dividendes de $\delta$ calcule une fair value trop élevée de $\delta e^{rT}$, voit le
contrat à terme comme bon marché de ce montant et l'achète~: l'arbitrage apparent n'est autre que l'erreur de prévision, et il perd lorsque les dividendes sont versés.
\end{solution}
